\section{Frequency tables and empirical probabilities}

For discrete or grouped data, the first natural summary is a frequency table. A
frequency table records how often each value, or each class of values, occurs in
the dataset.

Suppose the observed values are
\[
2,\quad 3,\quad 2,\quad 5,\quad 3,\quad 3,\quad 4,\quad 2.
\]
There are \(n=8\) observations. The value \(2\) appears three times, the value
\(3\) appears three times, the value \(4\) appears once, and the value \(5\)
appears once.

\begin{center}
\begin{tabular}{ccc}
\toprule
Value \(a\) & Count \(n_a\) & Relative frequency \(n_a/n\) \\
\midrule
2 & 3 & \(3/8\) \\
3 & 3 & \(3/8\) \\
4 & 1 & \(1/8\) \\
5 & 1 & \(1/8\) \\
\bottomrule
\end{tabular}
\end{center}

\begin{definitionbox}
For a dataset \(x_1,\ldots,x_n\), the \textbf{frequency} of a value \(a\) is
\[
n_a=\#\{i:x_i=a\}.
\]
The \textbf{relative frequency} of \(a\) is
\[
\frac{n_a}{n}.
\]
\end{definitionbox}

The relative frequencies add to \(1\), because every observation contributes to
exactly one observed value:
\[
\sum_a \frac{n_a}{n}=1.
\]
This makes relative frequencies look like probabilities. Indeed, they define an
empirical probability law on the observed values.

\subsection*{Empirical probability}

The empirical probability of observing the value \(a\) in the dataset is the
relative frequency
\[
P_n(\{a\})=\frac{n_a}{n}.
\]
The notation \(P_n\) emphasizes that this probability is built from a sample of
size \(n\). It is not necessarily the true probability from the population or
from the theoretical model.

For the example above,
\[
P_n(\{2\})=\frac38,
\qquad
P_n(\{4\})=\frac18.
\]
These are exact statements about the observed data. They do not prove that the
true probability of observing \(2\) is \(3/8\). They say that \(3/8\) of this
particular sample consists of the value \(2\).

\begin{remarkbox}
Relative frequency is data-based. Probability in a model is theory-based. The
relationship between them is one of the central themes of statistics, but the two
objects should not be identified without explanation.
\end{remarkbox}

\subsection*{Grouped frequencies}

For continuous data, exact repeated values may be rare. If waiting times are
recorded with many decimal places, it may happen that no two observations are
exactly equal. A table of exact values is then not very informative.

Instead, values are often grouped into intervals. For example, suppose waiting
times are grouped into bins
\[
[0,2),\quad [2,4),\quad [4,6),\quad [6,8).
\]
For each bin, we count how many observations fall inside it:
\[
n_j=\#\{i:a_j\leq x_i<a_{j+1}\}.
\]
The relative frequency of the bin is
\[
\frac{n_j}{n}.
\]
This is the empirical proportion of observations in that interval.

\subsection*{Connection with probability laws}

If a theoretical model says that a random variable \(X\) has probability law
\(P\), then the model probability of a set \(A\) is
\[
P(X\in A).
\]
For a dataset, the empirical counterpart is
\[
P_n(A)=\frac{\#\{i:x_i\in A\}}{n}.
\]
Thus, for an interval \([a,b]\),
\[
P_n([a,b])=\frac{\#\{i:a\leq x_i\leq b\}}{n}.
\]
This number tells us what proportion of the observed data lies in the interval.
It is the data-based analogue of the theoretical probability
\[
P(a\leq X\leq b).
\]

\subsection*{A warning about small samples}

In a small sample, relative frequencies can be unstable. If a coin is tossed ten
times and heads appears seven times, the relative frequency of heads is \(0.7\).
This does not by itself prove that the probability of heads is \(0.7\). It says
that in this sample, heads appeared in \(70\%\) of the trials.

As sample size increases, relative frequencies often stabilize under appropriate
probabilistic assumptions. The mathematical explanation of that stabilization
belongs to later chapters on sampling and limit theorems. Here we use relative
frequencies descriptively.

\begin{summarybox}
A frequency counts observations. A relative frequency divides that count by the
sample size. Relative frequencies form an empirical probability law on the data,
but they should not be confused with the unknown probabilities of the population
or model.
\end{summarybox}
